\documentclass[10pt,a4paper]{amsart}
\usepackage[margin=19mm]{geometry}
\usepackage{amsmath,amssymb,mathtools}
\usepackage[scr=rsfs,cal=euler]{mathalfa}
\usepackage{xfrac}
\usepackage[unicode,hidelinks,bookmarks=false]{hyperref}
\newcommand{\ie}{i.e.\ }
\newcommand{\cf}{cf.\ }
\newcommand{\resp}{respectively\ }
\newcommand{\Subsection}{\subsection}
\providecommand{\nobreakdash}{\nobreak\hbox{-}}
\setlength{\emergencystretch}{2em}
\multlinegap0pt
\usepackage{bm}
\usepackage{bbm}
\usepackage{dsfont}
\usepackage{xspace}
\usepackage{cancel}
\usepackage{mathtools}
\usepackage{amsthm}
\makeatletter
\@ifundefined{theorem}{\newtheorem{theorem}{Theorem}}{}
\@ifundefined{proposition}{\newtheorem{proposition}[theorem]{Proposition}}{}
\makeatother
\makeatletter
\newcommand{\Folner}{F\o{}lner}
\def\N{{\mathbb N}}
\def\Z{{\mathbb Z}}
\expandafter\ifx\csname customthm\endcsname\relax
\newenvironment{customthm}[1]
{\renewcommand\theinnercustomthm{#1}\innercustomthm}
\fi
\newcommand*\diff{\mathop{}\!\mathrm{d}}
\newcommand{\oh}{{\rm o}}
\makeatother
\title[Sharp density conditions for infinite $B+B$ sumsets in abeli]{Sharp density conditions for infinite $B+B$ sumsets in abelian groups}
\author{Ioannis Kousek}
\date{Source: arXiv:2607.18132v2 (CC BY 4.0)}
\begin{document}
\maketitle

\noindent\emph{Source and licence.} Sharp density conditions for infinite $B+B$ sumsets in abelian groups. Version arXiv:2607.18132v2 (CC BY 4.0), distributed under the Creative Commons Attribution 4.0 International licence, as recorded in the arXiv licence metadata for this exact version and shown on the abstract page https://arxiv.org/abs/2607.18132v2. Full rights notice: licenses/arxiv-2607.18132-CC-BY-4.0.txt.\par\smallskip

\noindent\textit{Editorial scope.} Counted unit: the Proposition whose source label is \texttt{no reverse implication of main theorems} in the article (source0.tex lines 943--950, source label \texttt{no reverse implication of main theorems}), delivered together with the article's own complete original proof of it (source0.tex lines 952--1025, one \texttt{proof} environment of 74 lines). No mathematics is written, repaired or re-derived: statement and proof are the article's own text line for line. Required context retained uncounted: main theorem (lines 308--325). Every definition, display and label of those lines is retained unchanged. Omitted from this package and not redistributed: 1--307, 326--942, 951--951, 1026--1250. Nothing inside a retained excerpt is omitted or altered: every retained body is delivered exactly as the article's lines are written, so no omission receipt of any kind accompanies this package. Citations: the retained excerpts carry no citation command at all, so no bibliography entry of the article is transcribed and no reference is redistributed. Reference adaptation: none was needed and none was made; every reference inside the delivered unit resolves inside it, so no unresolved reference or citation is produced. Printed slips kept literal: no printed word, symbol, hypothesis or number of the counted statement or of its proof was altered. Layout and class: the article's own class is \texttt{amsart} with packages including fontenc, pgf, pgfarrows, pgfnodes, pgfautomata, pgfheaps, pgfshade, hyperref, amssymb, babel, cleveref, mathtools, todonotes, soul; the wrapper is the lane's pinned \texttt{amsart} editorial wrapper at 10pt on A4, which restates the theorem-like environments the retained text needs and adds the article's own macro definitions that the retained text needs (\texttt{\textbackslash Folner}, \texttt{\textbackslash N}, \texttt{\textbackslash Z}, \texttt{\textbackslash diff}, \texttt{\textbackslash oh}), with extra wrapper packages (bm, bbm, dsfont, xspace, cancel, mathtools). The delivered numbering is the unit's own. Assets: no retained span contains a figure environment, a graphics-inclusion command, a PDF-page inclusion, a table or a TikZ picture; no image file is copied into this package, the package declares no asset dependency and no image file is redistributed with it. Licence: version arXiv:2607.18132v2 of this article is distributed under the Creative Commons Attribution 4.0 International licence, as recorded in the arXiv licence metadata for this exact version and shown on the abstract page https://arxiv.org/abs/2607.18132v2. A full rights notice is delivered at licenses/arxiv-2607.18132-CC-BY-4.0.txt. Characters: every retained excerpt is pure ASCII, so the delivered engine, XeLaTeX with its default Latin Modern text font, reports no missing character.
\begin{theorem}\label{main theorem}
Let $(G,+)$ be a countable abelian group with $\ell=[G:2G]<\infty$ and $r=|\ker(D)|<\infty$. Let $A\subset G$ and $\Phi$ be any \Folner{} sequence in $G$ such that the densities below exist. Then, the following hold:
\begin{enumerate}
\item  \label{main_theorem_1_2}
If 
$\diff_\Phi(A)+ \diff_{\Phi/2}(A)> 1,$
there exist an infinite set $B\subset G$ and some $t\in G$ such that 
$t+B+B\subset A$. 
\item  \label{main_theorem_1_1} 
\vspace{2mm} 
If 
$\diff_\Phi(A) + \diff_{\Phi/2}(A)> \frac{2\ell-1}{\ell},$
there exists an infinite set $B\subset G$ such that $B+B\subset A$. 
\item   \label{main_theorem_evens}  \vspace{2mm} If 
$ \diff_\Phi(A\cap 2G)+\diff_{\Phi/2}(A\cap 2G) > \frac{1}{\ell},$
there exists an infinite set $B\subset G$ such that $B+B\subset A$.
\end{enumerate}    
\end{theorem}

\begin{proposition}\label{no reverse implication of main theorems}
There is a set 
$A\subset \Z$ such that 
$$\overline{\diff}_{\Phi}(A) \leq 1-\frac{\ell \alpha_{\Phi}}{\ell+r},$$
for every \Folner\ sequence $\Phi$ in $\Z$, but 
$$\diff_{F}(A)+\diff_{F/2}(A)>1, $$
for some \Folner\ sequence $F$ in $\Z$.
\end{proposition}

\begin{proof}
We let $F=(F_N)_{N\in \N}$ be a \Folner\ sequence that is 
not q.i.d. and such that $F_1,F_1/2,F_1/4,F_2,F_2/2,F_2/4,\ldots$ are pairwise disjoint. For example, let $F_N=[4^{2N},2\cdot 4^{2N})\cap \N$, for 
$N\in \N$. In particular, $F_N/2=[2\cdot 4^{2N-1},4^{2N})\cap \N$ and $F_N/4=[4^{2N-1},2\cdot 4^{2N-1})\cap \N$. We 
then consider a set $A\subset \N$ defined via 
$$A= \left(\bigcup_{N=1}^{\infty} F_N\right) \bigcup \left( \bigcup_{N=1}^{\infty} F_N/2 \cap k\N \right),$$
with $k\in \N$, $k\geq 3$. Essentially, this could 
be generalized by replacing $k\N$ with a syndetic set of 
sufficiently small density. 

First, we observe that 
$$\diff_{F}(A)+\diff_{F/2}(A)=1+\frac{1}{k}>1.$$ 
We remark that, in view of Theorem \ref{main theorem}, this 
implies the existence of some infinite set $B\subset \N$ and 
$t\in \N_0$ such that $t+B+B \subset A$; this is, however, 
not a necessary part of this proof. 

We proceed to prove that for every \Folner\ sequence $\Phi$ 
in 
$\Z$, it holds that 
\begin{equation}\label{eq:4} 
\overline{\diff}_{\Phi}(A) \leq 1-\frac{\ell \alpha_{\Phi}}{\ell+r}.    
\end{equation}
By the definition of $A$, we only have to consider \Folner\ 
sequences that only intersect non-trivially with $F$ and 
$F/2$. For, if $\Phi_N= G_N \cup B_N$, where $G_N \subset \bigcup_{N=1}^{\infty} F_{N}\cup F_{N}/2$ and $B_N\cap \left(\bigcup_{N=1}^{\infty} F_{N}\cup F_{N}/2\right)=\emptyset$, it holds that
$$\overline{\diff}_{\Phi}(A)=\limsup_{N\to \infty} \frac{|A\cap \Phi_N|}{|\Phi_N|}= \limsup_{N\to \infty} \frac{|A\cap G_N|}{|G_N|+|B_N|} \leq \limsup_{N\to \infty} \frac{|A\cap G_N|}{|G_N|}.$$
% we also need to check that if 
% $$\limsup_{N\to \infty} \frac{|A\cap G_N|}{|G_N|}\leq 1-\frac{2\alpha_G}{3}, then $$\limsup_{N\to \infty} \frac{|A\cap \Phi_N|}{|\Phi_N|} \leq  1-\frac{2\alpha_{\Phi}}{3}$$, but this works by the above and writing 
%$$\limsup_{N\to \infty} \frac{|A\cap G_N|}{|G_N|+|B_N|}=
%\limsup_{N\to \infty} \frac{|A\cap G_N|}{|G_N|} \cdot \frac{|G_N|}{|G_N|+|B_N|}$$
%and then also comparing the $\alpha$s by opening up the definition. 
To easy the presentation, we first address two simple cases.
If $\Phi=F$ or $\Phi=F/2$, we have that $\alpha_{F}=0$, and 
so \eqref{eq:4} holds by default. Next, we consider the case that $(\Phi_N)_{N\in \N}=(F_N\cup F_N/2)_{N\in \N}$. Then we have 
that 
$$\alpha_{\Phi}= \lim_{N\to \infty} 
\frac{|F_N/2|}{|F_N|+|F_N/2|}=\frac{1}{3}.$$
In the integers we have that $\ell=2$ and $r=1$, so this means 
that the right hand side on \eqref{eq:4} equals $7/9$. At the 
same time 
$$\overline{\diff}_{\Phi}(A)=\lim_{N\to \infty} \frac{|F_N|+|F_N/2|/k}{|F_N|+|F_N/2|} = \frac{2+1/k}{3}=\frac{2k+1}{3k} \leq \frac{7}{9}, $$
for $k\geq 3$.
Similar calculations show \eqref{eq:4} in the case that 
$\Phi_N=(G_N \cup H_N)_{N\in \N}$, where $G_N\subset F_N$ and $H_N\subset F_N/2$.

Finally, we consider the general case where $\Phi_N=(G_N \cup H_N)_{N\in \N}$, with $G_N=G_{N_{k_1}}\cup \cdots \cup G_{N_{k_N}}$ for $G_{N_{k_i}}\subset F_{n(N,k_i)}$, and $H_N=H_{N_{m_1}}\cup \cdots \cup H_{N_{m_N}}$ for 
$H_{N_{m_i}} \subset F_{m(N,m_i)}/2$. Note that, $(G_N)_{N\in \N}$ and $(H_N)_{N\in \N}$ have to be \Folner\ sequences also. 
In order to verify \eqref{eq:4} it 
suffices to show that for every $N\in \N$, 
$$\left| A\cap \Phi_N \right| + \frac{2}{3} \left| \Phi_N\cap \Phi_N/2 \right| \leq \left| \Phi_N \right| + \oh_{N\to\infty}(1),$$
which is implied by
\begin{equation}\label{eq:5}
3\left| A\cap \Phi_N \right| + 2\left| \Phi_N\cap \Phi_N/2 \right| \leq 3 \left| \Phi_N \right| + \oh_{N\to\infty}(1).  
\end{equation}
Now, using the definition of $\Phi_N$ and $A$, we see that the 
left hand side in \eqref{eq:5} equals 
$$ 3|G_N|+\frac{3}{k}|H_N|+2 |G_N/2\cap H_N| + \oh_{N\to\infty}(1),$$
which equals
$$3\sum_{i=1}^{k_N}\left|G_{N_{k_i}}\right| + \frac{3}{k}\sum_{j=1}^{m_N} \left| H_{N_{m_j}} \right| + 
2 \sum_{i,j} \left| G_{N_{k_i}}/2 \cap H_{N_{m_j}}  \right| + \oh_{N\to\infty}(1). $$
Note that we have here used the fact that $(H_N)$ is \Folner\ and the assumption that 
$\{F_N,F_N/2,F_N/4:N\in \N\}$ are all disjoint, because otherwise 
we could have other non-trivial intersections in 
$$\Phi_N\cap \Phi_N/2=(G_N\cup H_N)\cap (G_N/2\cup H_N/2).$$ 
But observe that $G_{N_{k_i}}/2 \cap H_{N_{m_j}} \neq \emptyset$ only if $F_{n(N,k_i)}/2 \cap F_{m(N,m_i)}/2 \neq \emptyset$, i.e. only if $n(N,k_i)=m(N,m_i)$. Note also that in this case, we have $|G_{N_{k_i}}/2 \cap H_{N_{m_j}}|\leq |H_{N_{m_j}}|$. 

For $k\geq 3$ it thus follows that
\begin{align*}
& 3\left| A\cap \Phi_N \right| + 2\left| \Phi_N\cap \Phi_N/2 \right| \leq 3 \sum_{i=1}^{k_N}\left|G_{N_{k_i}}\right| + 3 \sum_{j=1}^{m_N}\left|H_{N_{k_j}}\right| + \oh_{N\to\infty}(1) \\ 
=& 3|G_N|+3|H_N|+\oh_{N\to\infty}(1)=3|\Phi_N|+\oh_{N\to\infty}(1),    
\end{align*}
showing \eqref{eq:5} and therefore the result. 
\end{proof}

\end{document}
